\section{Class numbers and the cyclotomic enlargement}
\label{sl-section}

We prove two of the inputs isolated in Section~\ref{sc-section}:
the class-number estimate and the cyclotomic enlargement
assertion~\ref{sc-layer}.  The latter includes the inequality for
strict Selmer lengths.  It is a statement about this particular
cyclotomic tower and nonzero twist; it is not an assertion about
restriction in an arbitrary extension of prime-power degree.
The finite-field homology construction and the comparison with
units in assertions~\ref{sc-few} and \ref{sc-count} are separate.

\subsection{An elementary bound for the class number}

For a number field $F$, write $n=[F:\mathbb Q]=r_1+2r_2$,
$D=|\operatorname{disc}F|$, and let $N\mathfrak a$ be the
absolute norm of a nonzero integral ideal.  Use ordinary Lebesgue
measure on $\mathbb R^{r_1}\times\mathbb C^{r_2}$, with one
embedding from each conjugate pair.  The image of $\mathfrak a$
has covolume
\[
                 2^{-r_2}\sqrt D\,N\mathfrak a .
                 \tag{SL1}
\]
Indeed, for $\mathcal O_F$ the matrix of all embeddings has
squared determinant equal to the discriminant.  Replacing each
pair of complex rows by its real and imaginary parts changes
the absolute determinant by $1/2$.  Passing to an ideal
multiplies the covolume by its index $N\mathfrak a$.

\begin{lemma}\label{sl-small-ideal}
Every ideal class of $F$ contains a nonzero integral ideal of norm
at most
\[
          \left(\frac4\pi\right)^{r_2}
                 \frac{n!}{n^n}\sqrt D
             \ \le\ \sqrt D .
             \tag{SL2}
\]
\end{lemma}
\begin{proof}
The elementary lattice-volume argument in the proof of
Lemma~\ref{ac-positive-gram} shows that an open, convex,
centrally symmetric body of volume greater than $2^n$ times a
lattice's covolume contains a nonzero lattice point.
Apply it to the image of an integral ideal $\mathfrak a$ and
the body
\[
 B_T=\left\{(x,z):
       \sum_{i=1}^{r_1}|x_i|+
                  2\sum_{j=1}^{r_2}|z_j|<T\right\}.
\]
Integration in polar coordinates gives
\[
       \operatorname{vol}(B_T)
             =2^{r_1}\left(\frac\pi2\right)^{r_2}
                         \frac{T^n}{n!}.
                         \tag{SL3}
\]
For clarity, put $u_i=|x_i|$ and $v_j=2|z_j|$.
The measures become $2\,du_i$ and $(\pi/2)v_j\,dv_j$.
The integral of $\prod_jv_j$ over
$u_i,v_j\ge0$, $\sum u_i+\sum v_j<T$, is $T^n/n!$;
successive integration, using
$\int_0^T t^a(T-t)^b\,dt=a!b!T^{a+b+1}/(a+b+1)!$,
proves this formula.

Thus, whenever
$T^n>(4/\pi)^{r_2}n!\sqrt D\,N\mathfrak a$, there is
$0\ne\alpha\in\mathfrak a$ in $B_T$.  The arithmetic--geometric
mean inequality, applied to the $n$ nonnegative numbers
consisting of each $|x_i|$ and two copies of each $|z_j|$, gives
\[
        |N_{F/\mathbb Q}\alpha|
              \le(T/n)^n .
\]
Let $T$ decrease to the indicated threshold.  A bounded region
contains only finitely many lattice points, so one nonzero
$\alpha$ attains the resulting non-strict bound.
The ideal $(\alpha)\mathfrak a^{-1}$ is integral and has norm
$|N_{F/\mathbb Q}\alpha|/N\mathfrak a$.  Taking $\mathfrak a$
in the inverse of the desired class proves the first inequality.

For the second, induction gives $n!/n^n\le2^{1-n}$:
the ratio of the left-hand quantities at $n+1$ and $n$ is
$(n/(n+1))^n\le1/2$, by the binomial inequality
$(1+1/n)^n\ge2$.  Since $4/\pi<2$ and
$r_2\le n-1$ for $n\ge2$, the desired product is at most
$2^{1-n+r_2}\le1$.  For $n=1$ it is exactly one.
\end{proof}

\begin{lemma}\label{sl-ideal-count}
For every real $X\ge1$ the number of nonzero integral ideals
of norm at most $X$ is at most
\[
                    X(1+\log X)^{n-1}.
                    \tag{SL4}
\]
\end{lemma}
\begin{proof}
There are at most $n$ prime ideals above any rational prime $p$,
because their ramification indices and residue degrees satisfy
$\sum_{\mathfrak p\mid p}e_{\mathfrak p}f_{\mathfrak p}=n$.
For each $p$, assign distinct labels in $\{1,\ldots,n\}$ to
these prime ideals.  If
$\mathfrak a=\prod_{\mathfrak p}\mathfrak p^{a_{\mathfrak p}}$,
define integers $b_1,\ldots,b_n\ge1$ by placing the factor
$p^{a_{\mathfrak p}}$ in the coordinate with the label of
$\mathfrak p$.  Unique factorization of integers and ideals
shows that $\mathfrak a\mapsto(b_1,\ldots,b_n)$ is injective.
Moreover
\[
        \prod_i b_i
          =\prod_{\mathfrak p}p^{a_{\mathfrak p}}
          \le \prod_{\mathfrak p}
                         p^{f_{\mathfrak p}a_{\mathfrak p}}
          =N\mathfrak a.
\]
No assumption that the prime ideals are principal is involved.

Let $A_j(X)$ count $j$-tuples of positive integers with product
at most $X$.  The bound
$A_j(X)\le X(1+\log X)^{j-1}$ follows by induction.
For $j=1$, it is $\lfloor X\rfloor\le X$.  For $j\ge2$,
\[
 \begin{split}
 A_j(X)&=\sum_{a\le X}A_{j-1}(X/a)\\
       &\le X(1+\log X)^{j-2}\sum_{a\le X}\frac1a
        \le X(1+\log X)^{j-1}.
 \end{split}
\]
The last inequality follows by comparing the terms from $a=2$
onwards with the integral of $1/t$ over $[1,X]$.
The injection proves (SL4).
\end{proof}

\begin{corollary}\label{sl-class-number}
For every number field,
\[
       \log h_F\le\frac12\log D+
             (n-1)\log\left(1+\frac12\log D\right).
             \tag{SL5}
\]
In particular (SC14) and (SC15) hold.
\end{corollary}
\begin{proof}
Choose one ideal of norm at most $\sqrt D$ in each class by
Lemma~\ref{sl-small-ideal}, apply Lemma~\ref{sl-ideal-count}
with $X=\sqrt D$, and take logarithms.
The discriminant is a nonzero integer, so $D\ge1$ and all
quantities are defined.  Replacing $n-1$ by $n$ gives (SC14);
the calculation following that formula gives (SC15).
\end{proof}

\subsection{The cyclotomic degree and discriminant}

Let $p$ be prime and let $L$ contain $\mu_p$ if $p$ is odd,
and $\mu_4$ if $p=2$.  Write $c_0$ for the largest integer
such that $\mu_{p^{c_0}}\subset L$.  Then $c_0\ge1$ and
$c_0\ge2$ at $p=2$.  The cyclotomic character identifies
\[
 \operatorname{Gal}(L(\mu_{p^\infty})/L)
                         =1+p^{c_0}\mathbb Z_p,
 \qquad
 L_m=L(\mu_{p^{c_0+m}}),\quad [L_m:L]=p^m .
 \tag{SL6}
\]
Here the equality denotes the image under that character.
We recall the justification to fix both the degree and the
largest-root convention.

The polynomial $\Phi_{p^a}(1+X)$ is Eisenstein at $p$:
its constant term is $p$, its leading coefficient is one, and
its reduction modulo $p$ is $X^{p^{a-1}(p-1)}$.
For the last assertion use
$\Phi_{p^a}(Y)=\sum_{j=0}^{p-1}Y^{jp^{a-1}}$ and
the binomial identity in characteristic $p$.
Eisenstein's criterion follows directly by reducing a proposed
factorization into monic integral polynomials modulo $p$:
both constant terms would be divisible by $p$, contradicting
the constant term $p$.  The usual clearing of denominators
for primitive polynomials passes from rational to integral
factorizations.  Therefore
$[\mathbb Q(\mu_{p^a}):\mathbb Q]=p^{a-1}(p-1)$.
The automorphisms $\zeta\mapsto\zeta^b$ consequently fill the
whole unit group modulo $p^a$.

For $a\ge1$ (and $a\ge2$ at $2$), the group
$1+p^a\mathbb Z_p$ is topologically generated by $1+p^a$.
Indeed the valuation formula (SC18) shows that this element
has order $p^{b-a}$ modulo $p^b$, equal to the number of
elements in the indicated congruence group.  Every nontrivial
closed subgroup is $1+p^b\mathbb Z_p$ for some $b\ge a$:
choose a nonidentity element with least valuation of its
difference from one; the same valuation argument shows that
its powers fill that congruence group, and minimality puts
every element of the subgroup in it.
The image of the cyclotomic character for $L$ is a closed
subgroup of finite index in the cyclotomic Galois group.
It lies in $1+p\mathbb Z_p$, or in $1+4\mathbb Z_2$.
The largest-root definition forces its index parameter to
be $c_0$, proving (SL6).  It follows as well that the largest
$p$-power roots in $L_m$ have order exactly $p^{c_0+m}$.

\begin{lemma}\label{sl-cyclotomic-discriminant}
If $L'=L(\zeta)$, where $\zeta$ is a root of unity of order
$N\ge1$, and $d=[L':L]$, then
\[
                  D_{L'}\le D_L^d N^{[L':\mathbb Q]} .
                  \tag{SL7}
\]
\end{lemma}
\begin{proof}
Put $R=\mathcal O_L$, $B=\mathcal O_{L'}$ and
$A=R[\zeta]\subset B$.  The minimal polynomial $f$ of
$\zeta$ is monic in $R[X]$ and divides $X^N-1$ there.
The trace dual $A^\#=\{x:\operatorname{Tr}_{L'/L}(xA)
\subset R\}$ equals $f'(\zeta)^{-1}A$.
To see this, use the basis $1,\zeta,\ldots,\zeta^{d-1}$.
Lagrange interpolation on the distinct roots of $f$ gives
\[
 \operatorname{Tr}\bigl(\zeta^j/f'(\zeta)\bigr)
     =0\quad(0\le j<d-1),\qquad
 \operatorname{Tr}\bigl(\zeta^{d-1}/f'(\zeta)\bigr)=1.
\]
For larger $j$ these traces lie in $R$, by reducing $X^j$
modulo $f$.  Thus the pairing matrix for the bases
$\zeta^i$ and $\zeta^j/f'(\zeta)$ has zeros below the
opposite diagonal and ones on that diagonal.  Its determinant
is a unit in $R$.  This proves the trace-dual assertion over
$R$ itself, without assuming that $R$ is principal.

As $B^\#\subset A^\#$, we have
$f'(\zeta)B^\#\subset A\subset B$.
The relative different is
$\mathfrak D_{L'/L}=(B^\#)^{-1}$, so
$f'(\zeta)\in\mathfrak D_{L'/L}$.
Differentiating $X^N-1=f(X)g(X)$ and evaluating at $\zeta$
gives
$N\zeta^{N-1}=f'(\zeta)g(\zeta)$.
Since $\zeta$ is a unit, $N$ belongs to the different.
Hence $(N)\subset\mathfrak D_{L'/L}$ and
$N_{L'/\mathbb Q}\mathfrak D_{L'/L}\le N^{[L':\mathbb Q]}$.
The discriminant identity
$D_{L'}=D_L^d N_{L'/\mathbb Q}\mathfrak D_{L'/L}$
proves (SL7).  This identity follows by the determinant of
the trace pairing in a tower, or equivalently by transitivity
of the trace-dual fractional ideals.
\end{proof}

\subsection{The strict Selmer comparison in a layer}

\begin{proposition}\label{sl-selmer-monotonicity}
For $E\ge1$ and every $m\ge0$,
\[
                    s_p(L,E)\le s_p(L_m,E).
                    \tag{SL8}
\]
\end{proposition}
\begin{proof}
Use the common infinite tower and the compact module $X$
of Proposition~\ref{ct-descent}.  If $u$ is the cyclotomic
value of a generator and $T$ its action on $X$, put
$U=u^E T$.  That proposition identifies the two Pontryagin
duals with $X/(U-1)X$ and $X/(U^{p^m}-1)X$.
The polynomial identity
\[
 U^{p^m}-1=(U-1)(1+U+\cdots+U^{p^m-1})
\]
gives a surjection from the latter quotient to the former.
Dualizing gives an injection of the strict Selmer groups,
and the length inequality follows, with infinite length
allowed.  The descent proposition proves preservation of
the full local conditions; no assertion that ramification
conditions are automatically unchanged under an arbitrary
extension is being used.
\end{proof}

\begin{theorem}\label{sl-enlargement}
Assertion~\ref{sc-layer} is valid.
\end{theorem}
\begin{proof}
Let $n=[L:\mathbb Q]$ and choose the least $m\ge0$ such that
$p^m\ge4E+4$.  Put $d=p^m$, $L'=L_m$ and $c=c_0+m$.
Because $4E+4>1$, we have $m\ge1$, and minimality gives
\[
                    4E+4\le d<p(4E+4).
                    \tag{SL9}
\]
The hypotheses on roots of unity imply that $L$ is totally
imaginary and $n\ge2$.  They also imply $p\le n+1$:
for odd $p$ use $p-1=[\mathbb Q(\mu_p):\mathbb Q]\le n$;
for $p=2$ the inequality is immediate.
Consequently $d\le(n+1)(4E+4)$ and
$[L':\mathbb Q]=dn$.
The field $L'$ is totally imaginary, so
$r_2(L')=dn/2\ge d>2E+1$.
Also $p^c=p^{c_0}d\ge2d>2E+1$.

The largest-root assertion follows from (SL6).
Since $\mathbb Q(\mu_{p^c})\subset L'$,
\[
        p^{c-1}(p-1)\le n'=[L':\mathbb Q],
        \qquad p^c\le\frac{p}{p-1}n'\le2n'.
\]
Lemma~\ref{sl-cyclotomic-discriminant}, with $N=p^c$,
gives $D_{L'}\le D_L^d(p^c)^{n'}$.
Finally Proposition~\ref{sl-selmer-monotonicity} gives the
required length inequality.  These are all the assertions
in (SC13).
\end{proof}

Thus the exponent calculation in Proposition~\ref{sc-exponent}
no longer requires an unproved class-number estimate or
cyclotomic enlargement.  Its remaining assumptions are the
finite-field homology and class-group comparisons in
assertions~\ref{sc-few}--\ref{sc-sha}, to the extent not proved
in subsequent sections.


